\chapter{2006年福建卷（文）}
\section{单选题}
\begin{problem}
已知两条直线 \(y = ax - 2\) 和 \(y = (a + 2)x + 1\) 互相垂直，则 \(a\) 等于（\quad）

\choices
    {\(2\)}
    {\(1\)}
    {\(0\)}
    {\(-1\)}
\end{problem}

\begin{answer}
D
\end{answer}

\begin{solution}
两条直线的斜率分别为 \(a\) 和 \(a+2\). 由于两直线互相垂直，有
\(a(a+2)=-1\)，即 \(a^2+2a+1=0\)，所以 \((a+1)^2=0\)，从而 \(a=-1\)，故选 D.
\end{solution}

\begin{problem}
在等差数列 \(\{a_n\}\) 中，已知 \(a_1 = 2\)，\(a_2 + a_3 = 13\)，则 \(a_4 + a_5 + a_6\) 等于（\quad）

\choices
    {\(40\)}
    {\(42\)}
    {\(43\)}
    {\(45\)}
\end{problem}

\begin{answer}
B
\end{answer}

\begin{solution}
设等差数列的公差为 \(d\)，则 \(a_2=2+d\)，\(a_3=2+2d\). 由题意，
\(a_2+a_3=4+3d=13\)，解得 \(d=3\). 因此
\(a_4+a_5+a_6=(2+3d)+(2+4d)+(2+5d)=6+12d=42\)，故选 B.
\end{solution}

\begin{problem}
“\(\tan \alpha = 1\)”是“\(\alpha = \frac{\pi}{4}\)”的（\quad）

\choices
    {充分而不必要条件}
    {必要而不充分条件}
    {充要条件}
    {既不充分也不必要条件}
\end{problem}

\begin{answer}
B
\end{answer}

\begin{solution}
设命题 \(P\) 为“\(\tan\alpha=1\)”，命题 \(Q\) 为“\(\alpha=\frac{\pi}{4}\)”. 由 \(Q\) 可得 \(P\)，所以 \(P\) 是 \(Q\) 的必要条件. 但 \(P\) 成立时，\(\alpha=\frac{\pi}{4}+k\pi\)（\(k\in\Z\)），例如 \(\alpha=\frac{5\pi}{4}\) 时 \(P\) 成立而 \(Q\) 不成立，因此 \(P\) 不是 \(Q\) 的充分条件. 所以“\(\tan\alpha=1\)”是“\(\alpha=\frac{\pi}{4}\)”的必要而不充分条件，故选 B.
\end{solution}

\begin{problem}
已知 \(\alpha \in \left(\frac{\pi}{2}, \pi\right)\)，\(\sin \alpha = \frac{3}{5}\)，则 \(\tan \left(\alpha + \frac{\pi}{4}\right)\) 等于（\quad）

\choices
    {\(\frac{1}{7}\)}
    {\(7\)}
    {\(-\frac{1}{7}\)}
    {\(-7\)}
\end{problem}

\begin{answer}
A
\end{answer}

\begin{solution}
由 \(\alpha\in\left(\frac{\pi}{2},\pi\right)\) 可知 \(\cos\alpha<0\). 于是
\(\cos\alpha=-\sqrt{1-\sin^2\alpha}=-\frac{4}{5}\)，从而 \(\tan\alpha=-\frac{3}{4}\). 因为 \(1-\tan\alpha\neq0\)，所以
\[
\tan\left(\alpha+\frac{\pi}{4}\right)=\frac{\tan\alpha+1}{1-\tan\alpha}=\frac{-\frac{3}{4}+1}{1+\frac{3}{4}}=\frac{1}{7}
\]
故选 A.
\end{solution}

\begin{problem}
已知全集 \(U = \R\)，且 \(A = \{x \mid |x - 1| > 2\}\)，\(B = \{x \mid x^2 - 6x + 8 < 0\}\)，则 \((\complement_U A) \cap B\) 等于（\quad）

\choices
    {\([-1,4)\)}
    {\((2,3)\)}
    {\((2,3]\)}
    {\((-1,4)\)}
\end{problem}

\begin{answer}
C
\end{answer}

\begin{solution}
由 \(|x-1|>2\)，得 \(x<-1\) 或 \(x>3\)，所以 \(\complement_U A=[-1,3]\). 又
\(x^2-6x+8=(x-2)(x-4)\)，故 \(B=(2,4)\). 因此
\((\complement_U A)\cap B=[-1,3]\cap(2,4)=(2,3]\)，故选 C.
\end{solution}

\begin{problem}
函数 \(y = \dfrac{x}{x + 1}\)（\(x \neq -1\)）的反函数是（\quad）

\choices
    {\(y = \dfrac{x}{1 - x}\)（\(x \neq 1\)）}
    {\(y = \dfrac{x}{x - 1}\)（\(x \neq 1\)）}
    {\(y = \dfrac{x - 1}{x}\)（\(x \neq 0\)）}
    {\(y = \dfrac{1 - x}{x}\)（\(x \neq 0\)）}
\end{problem}

\begin{answer}
A
\end{answer}

\begin{solution}
由 \(y=\frac{x}{x+1}\) 得 \(y(x+1)=x\)，即 \(x(y-1)=-y\). 因为原函数的值域不含 \(1\)，所以 \(y\neq1\)，从而
\(x=\frac{y}{1-y}\). 交换 \(x\)，\(y\)，得反函数为 \(y=\frac{x}{1-x}\)，其定义域为 \(x\neq1\)，故选 A.
\end{solution}

\begin{problem}
已知正方体外接球的体积是 \(\frac{32}{3}\pi\)，那么正方体的棱长等于（\quad）

\choices
    {\(2\sqrt{2}\)}
    {\(\frac{2\sqrt{3}}{3}\)}
    {\(\frac{4\sqrt{2}}{3}\)}
    {\(\frac{4\sqrt{3}}{3}\)}
\end{problem}

\begin{answer}
D
\end{answer}

\begin{solution}
设正方体棱长为 \(s\)，则其外接球半径为体对角线的一半，即 \(R=\frac{\sqrt{3}}{2}s\). 由球的体积为 \(\frac{32}{3}\pi\)，有
\[
\frac{4}{3}\pi R^3=\frac{32}{3}\pi
\]
所以 \(R^3=8\)，即 \(R=2\). 因而 \(s=\frac{2R}{\sqrt{3}}=\frac{4\sqrt{3}}{3}\)，故选 D.
\end{solution}

\begin{problem}
从 \(4\) 名男生和 \(3\) 名女生中选出 \(3\) 人，分别从事三项不同的工作，若这 \(3\) 人中至少有 \(1\) 名女生，则选派方案共有（\quad）

\choices
    {\(108\) 种}
    {\(186\) 种}
    {\(216\) 种}
    {\(270\) 种}
\end{problem}

\begin{answer}
B
\end{answer}

\begin{solution}
从 \(7\) 人中选出 \(3\) 人并分别安排到三项不同工作，全部方案数为 \(\mathrm A_7^3=7\cdot6\cdot5=210\). 其中没有女生的方案是从 \(4\) 名男生中选派 \(3\) 人，方案数为 \(\mathrm A_4^3=4\cdot3\cdot2=24\). 所以至少有 \(1\) 名女生的方案数为
\(210-24=186\)，故选 B.
\end{solution}

\begin{problem}
已知向量 \(\bs{a}\) 与 \(\bs{b}\) 的夹角为 \(120^{\circ}\)，\(\left|\bs{a}\right| = 3\)，\(\left|\bs{a} + \bs{b}\right| = \sqrt{13}\)，则 \(\left|\bs{b}\right|\) 等于（\quad）

\choices
    {\(5\)}
    {\(4\)}
    {\(3\)}
    {\(1\)}
\end{problem}

\begin{answer}
B
\end{answer}

\begin{solution}
设 \(\left|\bs{b}\right|=t\)，由向量夹角公式，
\[
\left|\bs{a}+\bs{b}\right|^2=\left|\bs{a}\right|^2+t^2+2\left|\bs{a}\right|t\cos120^{\circ}=9+t^2-3t
\]
由 \(\left|\bs{a}+\bs{b}\right|=\sqrt{13}\)，得 \(t^2-3t-4=0\)，即 \((t-4)(t+1)=0\). 因为 \(t>0\)，所以 \(t=4\)，故选 B.
\end{solution}

\begin{problem}
对于平面 \(\alpha\) 和共面的直线 \(m\)，\(n\)，下列命题中真命题是（\quad）

\choices
    {若 \(m \perp \alpha\)，\(m \perp n\)，则 \(n \parallel \alpha\)}
    {若 \(m \parallel \alpha\)，\(n \parallel \alpha\)，则 \(m \parallel n\)}
    {若 \(m \subset \alpha\)，\(n \parallel \alpha\)，则 \(m \parallel n\)}
    {若 \(m\)，\(n\) 与 \(\alpha\) 所成的角相等，则 \(m \parallel n\)}
\end{problem}

\begin{answer}
C
\end{answer}

\begin{solution}
选项 C 中，\(m\subset\alpha\)，且 \(n\parallel\alpha\). 由于 \(m\)，\(n\) 共面，若 \(m\) 与 \(n\) 不平行，则两直线相交于某点，而交点属于 \(m\subset\alpha\)，这与 \(n\parallel\alpha\) 矛盾，所以 \(m\parallel n\)，选项 C 正确.

选项 B 中，两条直线都平行于同一平面时不一定互相平行；选项 D 中，与同一平面所成角相等也不推出两直线平行；选项 A 中，\(n\) 可能与平面 \(\alpha\) 相交，不能推出 \(n\parallel\alpha\). 因此真命题为 C.
\end{solution}

\begin{problem}
已知双曲线 \(\dfrac{x^2}{a^2} - \dfrac{y^2}{b^2} = 1\)（\(a > 0\)，\(b > 0\)）的右焦点为 \(F\)，若过点 \(F\) 且倾斜角为 \(60^{\circ}\) 的直线与双曲线的右支有且只有一个交点，则此双曲线离心率的取值范围是（\quad）

\choices
    {\(\left(1,2\right]\)}
    {\(\left(1,2\right)\)}
    {\(\left[2,+\infty\right)\)}
    {\(\left(2,+\infty\right)\)}
\end{problem}

\begin{answer}
C
\end{answer}

\begin{solution}
设双曲线的半焦距为 \(c\)，离心率为 \(e=\frac{c}{a}>1\)，则 \(c=ae\)，\(b^2=a^2(e^2-1)\). 过右焦点 \(F(c,0)\) 且倾斜角为 \(60^{\circ}\) 的直线为
\(y=\sqrt{3}(x-c)\). 令 \(t=\frac{x}{a}\)，将直线代入双曲线方程，得
\[
t^2-\frac{3(t-e)^2}{e^2-1}=1
\]
即
\[
(e^2-4)t^2+6et-4e^2+1=0
\]
将左边因式分解，得
\[
\bigl((e+2)t-(2e+1)\bigr)\bigl((e-2)t+(2e-1)\bigr)=0
\]
当 \(e\neq2\) 时，该方程的两个根为
\(t_1=\frac{2e+1}{e+2}\)，\(t_2=\frac{2e-1}{2-e}\).
双曲线右支上的点满足 \(x\geq a\)，即 \(t\geq1\). 对 \(e>1\)，始终有
\(t_1-1=\frac{e-1}{e+2}>0\)，所以 \(t_1\) 对应右支上的一个交点.

当 \(1<e<2\) 时，\(t_2=\frac{2e-1}{2-e}>1\)，所以有两个右支交点；当 \(e>2\) 时，\(t_2<0\)，所以只有一个右支交点. 当 \(e=2\) 时，代入上面的方程得 \(12t-15=0\)，也只有一个交点. 因此恰有一个交点时，离心率的取值范围为 \([2,+\infty)\)，故选 C.
\end{solution}

\begin{problem}
已知 \(f(x)\) 是周期为 \(2\) 的奇函数，当 \(0 < x < 1\) 时，\(f(x) = \lg x\). 设 \(a = f\left(\frac{6}{5}\right)\)，\(b = f\left(\frac{3}{2}\right)\)，\(c = f\left(\frac{5}{2}\right)\)，则（\quad）

\choices
    {\(a < b < c\)}
    {\(b < a < c\)}
    {\(c < b < a\)}
    {\(c < a < b\)}
\end{problem}

\begin{answer}
D
\end{answer}

\begin{solution}
利用周期性和奇偶性，
\[
a=f\left(\frac{6}{5}\right)=f\left(-\frac{4}{5}\right)=-f\left(\frac{4}{5}\right)=-\lg\frac{4}{5}=\lg\frac{5}{4}
\]
\[
b=f\left(\frac{3}{2}\right)=f\left(-\frac{1}{2}\right)=-f\left(\frac{1}{2}\right)=-\lg\frac{1}{2}=\lg2
\]
\[
c=f\left(\frac{5}{2}\right)=f\left(\frac{1}{2}\right)=\lg\frac{1}{2}=-\lg2
\]
因为 \(c<0<a\)，且 \(1<\frac{5}{4}<2\) 且对数函数 \(\lg x\) 在正数范围内单调递增，所以 \(c<a<b\)，故选 D.
\end{solution}

\section{填空题}
\begin{problem}
\(\left(x^{2} - \frac{1}{x}\right)^{5}\) 展开式中 \(x^{4}\) 的系数是\fillinblank{}.（用数字作答）
\end{problem}

\begin{answer}
10
\end{answer}

\begin{solution}
展开式的通项可写为从 \(5\) 个因式中选 \(k\) 个取 \(-\frac{1}{x}\)，其余因式取 \(x^2\)，所得项的幂为
\(x^{2(5-k)}x^{-k}=x^{10-3k}\). 要得到 \(x^4\)，需满足 \(10-3k=4\)，解得 \(k=2\). 因此 \(x^4\) 的系数为
\(\mathrm C_5^2(-1)^2=10\).
\end{solution}

\begin{problem}
已知直线 \(x - y - 1 = 0\) 与抛物线 \(y = ax^2\) 相切，则 \(a =\)\fillinblank{}.
\end{problem}

\begin{answer}
\(\frac{1}{4}\)
\end{answer}

\begin{solution}
直线方程给出 \(y=x-1\). 与抛物线 \(y=ax^2\) 联立，得
\(ax^2-x+1=0\). 直线与抛物线相切时，该方程有且只有一个实根，所以判别式满足
\(\Delta=(-1)^2-4a=0\)，从而 \(a=\frac{1}{4}\).
\end{solution}

\begin{problem}
已知实数 \(x\)，\(y\) 满足 \(\begin{cases}
y \leqslant 1,\\
y \geqslant |x - 1|.
\end{cases}\)，则 \(x + 2y\) 的最大值是\fillinblank{}.
\end{problem}

\begin{answer}
4
\end{answer}

\begin{solution}
由 \(y\geq|x-1|\geq0\) 及 \(y\leq1\)，得 \(0\leq y\leq1\). 对固定的 \(y\)，不等式 \(|x-1|\leq y\) 等价于
\(1-y\leq x\leq1+y\). 因为目标式 \(x+2y\) 随 \(x\) 增大而增大，所以在固定 \(y\) 时取 \(x=1+y\)，于是
\(x+2y\leq(1+y)+2y=1+3y\leq4\).
当 \(x=2\)，\(y=1\) 时等号成立，所以最大值为 \(4\).
\end{solution}

\begin{problem}
已知函数 \(f(x) = 2\sin \omega x\)（\(\omega > 0\)）在区间 \(\left[-\frac{\pi}{3}, \frac{\pi}{4}\right]\) 上的最小值是 \(-2\)，则 \(\omega\) 的最小值等于\fillinblank{}.
\end{problem}

\begin{answer}
\(\frac{3}{2}\)
\end{answer}

\begin{solution}
函数的最小值为 \(-2\)，而 \(2\sin\omega x\geq-2\)，所以区间内必须存在 \(x\)，使 \(\sin\omega x=-1\). 这要求
\(\omega x=-\frac{\pi}{2}+2k\pi\)（\(k\in\Z\)）.

当 \(0<\omega<\frac{3}{2}\) 时，区间 \(\left[-\frac{\pi}{3},\frac{\pi}{4}\right]\) 经乘以 \(\omega\) 后的左端点大于 \(-\frac{\pi}{2}\)，右端点小于 \(\frac{3\pi}{2}\)，因此不含任何 \(-\frac{\pi}{2}+2k\pi\)，函数不能取得 \(-2\). 当 \(\omega=\frac{3}{2}\) 时，取 \(x=-\frac{\pi}{3}\)，有 \(\omega x=-\frac{\pi}{2}\)，故函数值为 \(-2\). 因而 \(\omega\) 的最小值为 \(\frac{3}{2}\).
\end{solution}

\section{解答题}
\begin{problem}
已知函数 \(f(x) = \sin^{2}x + \sqrt{3}\sin x\cos x + 2\cos^{2}x\)，\(x \in \R\).

\begin{enumerate}
\item 求函数 \(f(x)\) 的最小正周期和单调增区间；
\item 函数 \(f(x)\) 的图象可以由函数 \(y = \sin 2x\)（\(x \in \R\)）的图象经过怎样的变换得到？
\end{enumerate}
\end{problem}

\begin{solution}
\begin{enumerate}
\item 利用二倍角公式，\[
f(x)=\frac{1-\cos2x}{2}+\frac{\sqrt{3}}{2}\sin2x+1+\cos2x
=\frac{3}{2}+\sin\left(2x+\frac{\pi}{6}\right).
\]
因此 \(f(x)\) 的最小正周期为 \(\pi\). 又
\(f'(x)=2\cos\left(2x+\frac{\pi}{6}\right)\)，所以 \(f'(x)>0\) 等价于
\[
-\frac{\pi}{2}+2k\pi<2x+\frac{\pi}{6}<\frac{\pi}{2}+2k\pi
\]
即 \(k\pi-\frac{\pi}{3}<x<k\pi+\frac{\pi}{6}\)，其中 \(k\in\Z\). 这就是函数的单调增区间.
\item 因为
\(f(x)=\sin\left(2\left(x+\frac{\pi}{12}\right)\right)+\frac{3}{2}\)，所以将 \(y=\sin2x\) 的图象向左平移 \(\frac{\pi}{12}\) 个单位长度，再向上平移 \(\frac{3}{2}\) 个单位长度即可得到 \(f(x)\) 的图象.
\end{enumerate}
\end{solution}

\begin{problem}
每次抛掷一枚骰子（六个面上分别标以数 \(1\)，\(2\)，\(3\)，\(4\)，\(5\)，\(6\)）.

\begin{enumerate}
\item 连续抛掷 \(2\) 次，求向上的数不同的概率；
\item 连续抛掷 \(2\) 次，求向上的数之和为 \(6\) 的概率；
\item 连续抛掷 \(5\) 次，求向上的数为奇数恰好出现 \(3\) 次的概率.
\end{enumerate}
\end{problem}

\begin{solution}
\begin{enumerate}
\item 连续抛掷 \(2\) 次的有序结果共有 \(6^2=36\) 种，其中两次数字相同的结果有 \(6\) 种，所以两次数字不同的概率为
\(\frac{36-6}{36}=\frac{5}{6}\).
\item 两次数字之和为 \(6\) 的有序结果为 \((1,5)\)，\((2,4)\)，\((3,3)\)，\((4,2)\)，\((5,1)\)，共 \(5\) 种，所以所求概率为 \(\frac{5}{36}\).
\item 每次向上的数为奇数的概率为 \(\frac{3}{6}=\frac{1}{2}\)，为偶数的概率也为 \(\frac{1}{2}\). 恰好出现 \(3\) 次奇数的概率为
\(\mathrm C_5^3\left(\frac{1}{2}\right)^3\left(\frac{1}{2}\right)^2=\frac{10}{32}=\frac{5}{16}\).
\end{enumerate}
\end{solution}

\begin{problem}
如图，四面体 \(ABCD\) 中，\(O\)，\(E\) 分别是 \(BD\)，\(BC\) 的中点，\(CA = CB = CD = BD = 2\)，\(AB = AD = \sqrt{2}\).
\begin{enumerate}
\item 求证：\(AO \perp\) 平面 \(BCD\)；
\item 求异面直线 \(AB\) 与 \(CD\) 所成角的大小；
\item 求点 \(E\) 到平面 \(ACD\) 的距离.
\end{enumerate}

\bitmapfigure[width=0.34\linewidth]{img/2006/fujian_liberal/q19_fig1.png}
\end{problem}

\begin{solution}
\begin{enumerate}
\item 因为 \(BCD\) 是边长为 \(2\) 的等边三角形，可建立空间直角坐标系，使
\(B=(-1,0,0)\)，\(D=(1,0,0)\)，\(C=(0,\sqrt{3},0)\)，\(O=(0,0,0)\). 设 \(A=(u,v,w)\). 由 \(AB=AD\)，有
\[
(u+1)^2+v^2+w^2=(u-1)^2+v^2+w^2
\]
所以 \(u=0\). 再由 \(AB=\sqrt{2}\)，得 \(v^2+w^2=1\)，由 \(AC=2\)，得
\((v-\sqrt{3})^2+w^2=4\). 两式相减得 \(3-2\sqrt{3}v=3\)，所以 \(v=0\)，于是 \(w^2=1\). 关于平面 \(BCD\) 对称时结论相同，取 \(A=(0,0,1)\). 此时平面 \(BCD\) 的方程为 \(z=0\)，而 \(AO\) 的方向向量平行于 \(z\) 轴，所以 \(AO\perp\) 平面 \(BCD\).
\item 取 \(\overrightarrow{AB}=(-1,0,-1)\)，\(\overrightarrow{CD}=(1,-\sqrt{3},0)\). 则
\(\overrightarrow{AB}\cdot\overrightarrow{CD}=-1\)，\(\left|\overrightarrow{AB}\right|=\sqrt{2}\)，\(\left|\overrightarrow{CD}\right|=2\). 异面直线所成的角取锐角，因此
\[
\cos\theta=\frac{\left|\overrightarrow{AB}\cdot\overrightarrow{CD}\right|}{\left|\overrightarrow{AB}\right|\left|\overrightarrow{CD}\right|}=\frac{1}{2\sqrt{2}}=\frac{\sqrt{2}}{4}
\]
从而 \(\theta=\arccos\frac{\sqrt{2}}{4}\).
\item \(E\) 是 \(BC\) 的中点，所以 \(E=\left(-\frac{1}{2},\frac{\sqrt{3}}{2},0\right)\). 平面 \(ACD\) 经过 \(A=(0,0,1)\)，\(C=(0,\sqrt{3},0)\)，\(D=(1,0,0)\)，其方程为
\[
x+\frac{y}{\sqrt{3}}+z=1
\]
点 \(E\) 到该平面的距离为
\[
\frac{\left|-\frac{1}{2}+\frac{\frac{\sqrt{3}}{2}}{\sqrt{3}}+0-1\right|}{\sqrt{1+\frac{1}{3}+1}}
=\frac{1}{\sqrt{\frac{7}{3}}}
=\frac{\sqrt{21}}{7}
\]
\end{enumerate}
\end{solution}

\begin{problem}
已知椭圆 \(\dfrac{x^2}{2} + y^2 = 1\) 的左焦点为 \(F\)，\(O\) 为坐标原点.
\begin{enumerate}
\item 求过点 \(O\)，\(F\)，并且与椭圆的左准线 \(l\) 相切的圆的方程；
\item 设过点 \(F\) 且不与坐标轴垂直的直线交椭圆于 \(A\)，\(B\) 两点，并且线段 \(AB\) 的中点在直线 \(x + y = 0\) 上，求直线 \(AB\) 的方程.
\end{enumerate}

\bitmapfigure[width=0.42\linewidth]{img/2006/fujian_liberal/q20_fig1.png}
\end{problem}

\begin{solution}
\begin{enumerate}
\item 由 \(\frac{x^2}{2}+y^2=1\) 可知椭圆的半长轴为 \(\sqrt{2}\)，半短轴为 \(1\)，半焦距为 \(c=1\)，所以左焦点 \(F=(-1,0)\)，左准线 \(l\) 的方程为 \(x=-\frac{a^2}{c}=-2\). 设所求圆的圆心为 \((h,k)\)，半径为 \(r\). 因为圆经过 \(O=(0,0)\) 和 \(F=(-1,0)\)，有
\[
h^2+k^2=(h+1)^2+k^2
\]
从而 \(h=-\frac{1}{2}\)，且 \(r^2=\frac{1}{4}+k^2\). 圆与直线 \(x=-2\) 相切，所以圆心到该直线的距离等于半径，即
\[
r=\left|-\frac{1}{2}+2\right|=\frac{3}{2}
\]
于是 \(\frac{1}{4}+k^2=\frac{9}{4}\)，解得 \(k=\pm\sqrt{2}\). 所以圆的方程为
\[
\left(x+\frac{1}{2}\right)^2+\left(y\pm\sqrt{2}\right)^2=\frac{9}{4}
\]
\item 设直线 \(AB\) 的斜率为 \(k\). 由于直线过 \(F=(-1,0)\) 且不与坐标轴垂直，\(k\neq0\)，其方程可写为 \(y=k(x+1)\). 代入椭圆方程，得
\[
(1+2k^2)x^2+4k^2x+2k^2-2=0
\]
设交点 \(A\)，\(B\) 的横坐标分别为 \(x_1\)，\(x_2\)，则其中点横坐标为
\[
x_M=\frac{x_1+x_2}{2}=-\frac{2k^2}{1+2k^2}
\]
又因 \(y_M=k(x_M+1)=\frac{k}{1+2k^2}\)，且中点在 \(x+y=0\) 上，所以
\[
x_M+y_M=\frac{k-2k^2}{1+2k^2}=0
\]
由 \(k\neq0\)，得 \(k=\frac{1}{2}\). 故直线 \(AB\) 的方程为
\[
y=\frac{1}{2}(x+1)
\]
即 \(x-2y+1=0\).
\end{enumerate}
\end{solution}

\begin{problem}
已知 \(f(x)\) 是二次函数，不等式 \(f(x) < 0\) 的解集是 \((0,5)\)，且 \(f(x)\) 在区间 \([-1,4]\) 上的最大值是 \(12\).

\begin{enumerate}
\item 求 \(f(x)\) 的解析式；
\item 是否存在自然数 \(m\)，使得方程 \(f(x) + \dfrac{37}{x} = 0\) 在区间 \((m,m+1)\) 内有且只有两个不等的实数根？若存在，求出所有 \(m\) 的值；若不存在，说明理由.
\end{enumerate}
\end{problem}

\begin{solution}
\begin{enumerate}
\item 因为 \(f(x)\) 是二次函数，且 \(f(x)<0\) 的解集为 \((0,5)\)，所以 \(0,5\) 是它的两个零点，且开口向上. 于是可设
\(f(x)=a x(x-5)\)，其中 \(a>0\). 在区间 \([-1,4]\) 上，\(f\) 为开口向上的二次函数，所以最大值在端点处取得. 又
\(f(-1)=6a\)，\(f(4)=-4a\).
由最大值为 \(12\)，得 \(6a=12\)，所以 \(a=2\)，从而
\(f(x)=2x(x-5)\).
\item 将 \(f(x)=2x(x-5)\) 代入方程，并注意区间 \((m,m+1)\) 中 \(x>0\)，可得
\[
2x(x-5)+\frac{37}{x}=0
\Longleftrightarrow
2x^3-10x^2+37=0
\]
设 \(g(x)=2x^3-10x^2+37\)，则
\(g'(x)=2x(3x-10)\). 所以 \(g\) 在 \((0,\frac{10}{3})\) 上递减，在 \((\frac{10}{3},+\infty)\) 上递增. 计算得
\(g(0)=37>0\)，\(g(3)=1>0\)，\(g\left(\frac{10}{3}\right)=-\frac{1}{27}<0\)，\(g(4)=5>0\).
因此 \(g(x)=0\) 在 \((3,\frac{10}{3})\) 和 \((\frac{10}{3},4)\) 内各有一个根，且这两个根不等，所以在 \((3,4)\) 内恰有两个不等的实数根. 在 \((0,3)\) 上 \(g(x)>0\)，在 \((4,+\infty)\) 上 \(g(x)>0\)，故其他自然数 \(m\) 均不满足条件. 所以存在且所有这样的 \(m\) 为 \(m=3\).
\end{enumerate}
\end{solution}

\begin{problem}
已知数列 \(\{a_n\}\) 满足 \(a_1 = 1\)，\(a_2 = 3\)，\(a_{n+2} = 3a_{n+1} - 2a_n\)（\(n \in \N^*\)）.

\begin{enumerate}
\item 证明：数列 \(\{a_{n+1} - a_n\}\) 是等比数列；
\item 求数列 \(\{a_n\}\) 的通项公式；
\item 若数列 \(\{b_n\}\) 满足 \(4^{b_1-1}4^{b_2-1}\cdots 4^{b_n-1} = (a_n + 1)^{b_n}\)（\(n \in \N^*\)），证明：\(\{b_n\}\) 是等差数列.
\end{enumerate}
\end{problem}

\begin{solution}
\begin{enumerate}
\item 令 \(d_n=a_{n+1}-a_n\). 由递推式，
\[
d_{n+1}=a_{n+2}-a_{n+1}=2a_{n+1}-2a_n=2d_n
\]
又 \(d_1=a_2-a_1=2\)，所以 \(\{d_n\}=\{a_{n+1}-a_n\}\) 是以 \(2\) 为首项、\(2\) 为公比的等比数列.
\item 由上问，\(a_{n+1}-a_n=2^n\). 因此
\[
a_n=a_1+\sum_{k=1}^{n-1}(a_{k+1}-a_k)
=1+\sum_{k=1}^{n-1}2^k
=1+(2^n-2)=2^n-1
\]
\item 由 \(a_n+1=2^n\)，题设等式化为
\[
4^{\sum_{k=1}^n(b_k-1)}=(2^n)^{b_n}
\]
由于 \(4=2^2\)，且指数函数 \(2^x\) 在实数范围内是单调函数，有
\[
2\sum_{k=1}^n(b_k-1)=nb_n
\]
当 \(n=1\) 时，得 \(2(b_1-1)=b_1\)，所以 \(b_1=2\). 对 \(n\geq2\)，上式在 \(n\) 和 \(n+1\) 时分别为
\(2\sum_{k=1}^n(b_k-1)=nb_n\)，\(2\sum_{k=1}^{n+1}(b_k-1)=(n+1)b_{n+1}\).
两式相减，得
\[
(n-1)b_{n+1}=nb_n-2
\]
于是
\[
\frac{b_{n+1}-2}{n}=\frac{b_n-2}{n-1}
\]
所以 \(\frac{b_n-2}{n-1}\)（\(n\geq2\)）为常数，设其为 \(q\)，则
\(b_n=2+(n-1)q\). 这正是等差数列的通项形式，故 \(\{b_n\}\) 是等差数列.
\end{enumerate}
\end{solution}
